\chapter*{Infinite Series}
\addcontentsline{toc}{chapter}{Infinite Series}

\section*{From Sequences to Series}

An infinite series is the sum of infinitely many terms:
$$\sum_{n=1}^{\infty} a_n = a_1+a_2+a_3+\cdots$$
But this sum is not interpreted directly. Instead, we study the sequence of \textbf{partial sums}
$$s_N = \sum_{n=1}^{N} a_n.$$ 

\begin{definition}[Convergence of a series]
The series $\sum a_n$ converges if the sequence of partial sums $\{s_N\}$ converges to a finite limit $S$:
$$\lim_{N\to\infty} s_N = S.$$ 
In that case, we write
$$\sum_{n=1}^{\infty} a_n = S.$$ 
Otherwise, the series diverges.
\end{definition}

\section*{Geometric Series}

The most important basic model is the geometric series
$$\sum_{n=0}^{\infty} ar^n = a+ar+ar^2+\cdots$$
with partial sums
$$s_N = a\frac{1-r^{N+1}}{1-r}, \qquad r\neq 1.$$ 

Taking limits gives:
\begin{theorem}[Geometric series test]
The geometric series converges if and only if $|r|<1$. In that case,
$$\sum_{n=0}^{\infty} ar^n = \frac{a}{1-r}.$$ 
If $|r|\ge 1$, the series diverges.
\end{theorem}

This explains Zeno-type sums such as
$$\frac12+\frac14+\frac18+\cdots = 1.$$ 

\section*{Telescoping Series}

A telescoping series is one in which large cancellations occur between consecutive partial sums. A standard pattern is
$$\sum_{n=1}^{\infty}(b_n-b_{n+1}).$$
Then
$$s_N = (b_1-b_2)+(b_2-b_3)+\cdots+(b_N-b_{N+1}) = b_1-b_{N+1}.$$ 

So the convergence of the series reduces to the limit of $b_{N+1}$.

\section*{Necessary Condition for Convergence}

\begin{theorem}[Test for divergence]
If the series $\sum a_n$ converges, then necessarily
$$\lim_{n\to\infty} a_n = 0.$$ 
\end{theorem}

\begin{remark}
This condition is necessary but not sufficient. For example,
$$\sum_{n=1}^{\infty} \frac{1}{n}$$
diverges even though $1/n\to 0$.
\end{remark}